Los resultados obtenidos en este estudio muestran que la teoría de la complejidad algorítmica 
se refleja de manera consistente en la práctica, aunque con algunas diferencias. 
En el caso de los algoritmos de ordenamiento, se confirmó que algoritmos $O(n \log n)$ como Merge Sort, Sort, y Quick Sort escalan de forma eficiente, solo que este último puede variar más según la distrubución y tamaño.
Esto se debe por la posición del pivote que genera aleatoriedad.  
Entre ellos, Sort se posiciona como la mejor alternativa, 
gracias a las optimizaciones de su implementación.

En la multiplicación de matrices, Strassen demuestra su superioridad teórica sobre el método Naive al reducir la complejidad a $O(n^{2.81})$. 
No obstante, en tamaños pequeños su sobrecarga recursiva hace que Naive sea más competitivo, lo que muestra que la eficiencia práctica 
depende no solo de la complejidad asintótica, sino también de la estructura del algoritmo y del contexto de ejecución.

A un nivel más general, este trabajo evidencia la relevancia del análisis experimental como complemento indispensable de la teoría en el aprendizaje de este curso. 
Las medidas de tiempo y memoria confirman que la elección del algoritmo adecuado no depende únicamente de su clasificación teórica, 
sino también del tamaño de la entrada, la naturaleza de los datos y la infraestructura disponible. 
De este modo, los hallazgos contribuyen a un mejor entendimiento de cómo teoría y práctica se entrelazan en la computación, 
y resaltan la necesidad de evaluar algoritmos en escenarios reales antes de aplicarlos.